\chapter*{初版前言}
\addcontentsline{toc}{chapter}{初版前言} % 手动加入目录
\markboth{初版前言}{初版前言} % 设置页眉

本集中的论文作于1962至1968年间，旨在重构黑格尔法哲学与政治的自然法的传统之间的联系，同时也重构向现代世界以及与之相应的工业革命和政治革命的那种转变。在时至今日仍然影响深远的法、国家和社会等概念的崭新规定中，黑格尔法哲学也引领了这种转向。这一目标的前提，殊异于流俗理解之下的《法哲学原理》的效果史及其对马克思主义、自由主义和纳粹主义这些全球性意识形态的诸般影响所提供的历史观。如果我们与时下主流观点相反，想要像作者理解他本人一样去理解他，由此出发进行解释的话，那么在我们对法哲学进行历史——考证的解释时，就会发现自己需要面对黑格尔的问题，这些问题是前人在黑格尔自己的时代意识和问题意识下向他提出来的。这里，我试图通过一种允许我们以柏拉图和亚里士多德以及霍布斯、卢梭和康德的眼光去重新阅读黑格尔法哲学的视角，对其被歪曲的效果史进行校正。唯有经过此番校正，才能对自黑格尔、马克思以来决定了现代法哲学和国家哲学的那种概念变异与视角改换的影响做出评价。

1962年面世的《黑格尔〈法哲学原理〉中的传统与革命》一文最先提出了后来进一步研究的主线。此文是献给卡尔·洛维特(Karl Löwith)65岁寿辰的。本书的所有研究也都敬献给他。

\begin{flushright}
M. 里德尔\\
1969年7月\\
于海德堡
\end{flushright}